\scanpage{35}
\mspage{17}
% archive foliation: 34

sous $\mathbb{F}_2^J$, que je vais désigner par $P_{M,\Delta}$
\[
(34)\qquad F(\Delta_M)=
\underset{\text{ens des faces de }\Delta}{\underbrace{P_{M,\Delta}}}
\quad\text{torseur sous }\mathbb{F}_2^J\qquad(J=M_0^*)
\]
On a
\[
(35)\qquad
\mathbb{F}_2^J\simeq
\underset{M_0}{\underset{\simeq}{V_J(\mathbb{F}_2)}}
\times\underset{\mathrm{diag}}{\underbrace{\mathbb{F}_2}}
\]
et l'opération sur $P_{M,\Delta}$ de l'élément non nul \deleted{\unsure{$\mathbb{F}_2$}} du facteur $\mathbb{F}_2$ n'est autre que l'antipodisme. Ceci posé on a un isom. canonique
\[
(36)\qquad P_{M,\Delta}=F(\Delta_M)\xleftarrow{\sim}P_M/Z
\qquad
\underset{\text{centre de }GL(M)}{\underset{=}{Z=\{\pm\mathrm{id}_M\}}}
\]
provenant de l'homomorphisme évident
\[
(37)\qquad
\diagram{\begin{tikzcd}[column sep=large]
P_M=\{(\xi_i)_{i\in J}\in M^{*J}\mid\xi_i^0=i\ \forall i\in J,\ \sum\xi_i=0\}
\arrow[d,"{(\xi_i)\longmapsto\{(\dot\xi_i)\}}"]\\
F(\Delta_M)=
\{f\in\mathcal P_3(\Delta_M)\mid
\underset{\text{induit par }\Delta\to J}{\underbrace{f\to J}}\ \text{bijectif}\}
\end{tikzcd}}
\]
\[
\text{où }\dot\xi=\Lambda^*\xi=\pm\xi\in\Delta\ \text{pour }\xi\in M^*.
\]

Ceci nous montre à nouveau que $\mathrm{Aut}(\Delta_M)\simeq GL(M)'$, de façon précise on a les iso. de suites exactes
\[
(38)\qquad
\diagram{\begin{tikzcd}
1\arrow[r]&
\overset{M_0\times\mathbb{F}_2}{\overbrace{\mathrm{End}(M_0)/\mathbb{F}_2\cdot\mathrm{id}_{M_0}}}
\arrow[r]\arrow[d,"\sim"']&
GL(M)'\arrow[r]\arrow[d,"\sim"']&
\overset{\unsure{\mathcal{G}_J|\mathfrak S_J}}{\overbrace{GL_{\mathbb{F}_2}(M_0)}}\arrow[r]\arrow[d,"\sim"']&1\\
1\arrow[r]&
\underset{\simeq M_0\times\mathbb{F}_2}{\underbrace{\mathbb{F}_2^J}}\arrow[r]&
\mathrm{Aut}_{\mathrm{oct}}(\Delta_M)\arrow[r]&
\unsure{\mathcal{G}_J|\mathfrak S_J}\arrow[r]&1
\end{tikzcd}}
\]

\margin{\uline{NB} La donnée d'un octaèdre (non orienté) $\Delta$ équivaut à celle d'un couple $(I,\underline\mu)$, où $I$ est un ens. à quatre éléments (l'ens de ses \unsure{bifaces}, qui définit l'oct. \uline{orienté} correspondant) et où $\underline\mu=\underline\omega(\Delta)$ est un ens. à deux éléments. — \unsure{On verra} que (plus bas) $\unsure{\mathcal E_\Delta|I}$ \illegible{3} $\unsure{\simeq\underline\mu\wedge\underline\omega_I}$ de $\underline\mu$ également (cf plus bas).}

A un octaèdre $(\Delta,\underline a_\Delta)$ ou $(\Delta\longrightarrow J)$ sont associés trois ens. à deux éléments remarquables, correspondant aux trois caractères \inserted{non triviaux} de $\mathrm{Aut}_{\mathrm{oct}}(\Delta)$ i.e. aux trois hom. non nuls de
\deleted{\unsure{$\mathrm{Aut}_{\mathrm{oct}}(\Delta)$}}
\[
\mathrm{Aut}_{\mathrm{oct}}(\Delta)_{\mathrm{ab}}
\simeq(\mathrm{Aut}_{\mathrm{oct}}^+(\Delta)\times\{1,\underline a_\Delta\})_{\mathrm{ab}}
\simeq\mathbb{F}_2\times\mathbb{F}_2.
\]
Ce sont respectivement les ensembles
